\section{Méthodologie et Architecture Implémentée}

Notre approche ne s'est pas construite en une seule étape. Elle est le fruit d'itérations successives, partant d'un constat d'échec sur les méthodes classiques pour aboutir à une architecture hybride et agentique. 

Afin de fluidifier l'organisation, nous avons rapidement choisi d'utiliser une structure très modulaire dans laquelle des classes abstraites définissent les interfaces de chaque étape de la pipeline. Chaque nouvel outil ou amélioration est alors développé selon ces interfaces et peut donc être permuté avec un autre sans créer de conflits. Cela nous a permis de travailler de manière indépendante sur différentes étapes du processus, et de facilement comparer différentes solutions. Par exemple, nous avons implémenté une classe pour chaque LLM utilisé (Mistral, Llama...) héritant d'une même classe abstraite, que l'on pouvait donc utiliser dans le reste du projet sans se soucier des différences entre les API de chaque modèle. Cette manière de procéder se retrouve dans tout le reste de la pipeline.

\subsection{Préparation des données : Du code brut aux "Chunks" sémantiques}

La qualité d'un système RAG dépend avant tout de la qualité de ses données. Contrairement à du texte littéraire, le code source d'Envision possède une structure hiérarchique forte qui ne supporte pas un découpage arbitraire.

\subsubsection{Le problème du découpage standard}
Les méthodes classiques de \textit{chunking} (découpage par nombre fixe de tokens ou par séparateurs de paragraphes) s'avèrent destructrices pour du code.
\begin{itemize}
    \item Si un découpage coupe une fonction en deux, le modèle perd la signature de la fonction (ses arguments) ou son corps logique.
    \item Le contexte est souvent situé en début de fichier (déclaration des tables), information cruciale pour comprendre les lignes situées 500 lignes plus bas.
\end{itemize}

\subsubsection{Notre solution : Le Envision Chunker}
Nous avons développé un module de parsing spécifique (\texttt{EnvisionParser}) qui analyse la structure syntaxique des fichiers \texttt{.nvn}. Plutôt que de découper aveuglément, nous identifions des blocs logiques :
\begin{itemize}
    \item \textbf{Blocs de définition :} \texttt{table}, \texttt{expect}, \texttt{read}.
    \item \textbf{Blocs logiques :} Délimitations par les commentaires spéciaux \texttt{///} utilisés dans Envision.
\end{itemize}
Le \texttt{EnvisionChunker} regroupe ensuite ces segments en unités cohérentes, tout en veillant à ce que chaque \textit{chunk} conserve les métadonnées essentielles (nom du fichier source, bibliothèques importées) pour garantir sa compréhension isolée par le LLM. Si un bloc dépasse le seuil maximal de tokens autorisé, le \texttt{chunker} procède à sa division. Pour pallier la perte de contexte aux points de rupture, nous avons implémenté une stratégie de chevauchement de quelques lignes entre les segments adjacents, facilitant l'analyse du code situé à la frontière entre deux sous-chunks.

\subsubsection{La méthode RAPTOR}
Pour améliorer les résultats du RAG sur des questions générales, sans lien direct avec un morceau de code spécifique, nous avons implémenté la méthode RAPTOR (Retrieval-Augmented Generation with Tree-based Retrieval). Le processus commence par l'analyse des fichiers source au niveau de base, où les documents sont parsés, découpés en chunks et enrichis de résumés pré-calculés, avant d'être transformés en embeddings vectoriels pour faciliter les comparaisons de similarité. Ensuite, une construction récursive de l'arbre s'opère : les chunks sont regroupés en clusters à l'aide d'algorithmes de réduction de dimension (UMAP) et de modélisation probabiliste (GMM), puis chaque cluster fait l'objet d'un résumé consolidé, créant ainsi des niveaux hiérarchiques successifs jusqu'à atteindre une vue d'ensemble.

Lors du RAG, les résumés des clusters peuvent être renvoyés au même titre que les chunks issus des scripts d'origine, l'idée étant que ces résumés sont plus pertinents pour des questions générales. Les chemins des fichiers dont viennent les chunks du cluster sont également fournis au LLM, lui permettant de les explorer plus en détail lors d'une itération ultérieure.


% \subsection{Stratégie de Recherche Hybride : Le Routeur}

% L'une des contributions majeures de notre architecture est de ne pas se reposer uniquement sur la recherche vectorielle. Nous avons identifié deux typologies de questions radicalement différentes, nécessitant des outils distincts.

% \subsubsection{La limite du Vectoriel face au Lexical}
% La recherche vectorielle (via FAISS et des embeddings type \textit{all-MiniLM-L6-v2}) excelle pour comprendre des concepts (\textit{« Comment calculer le stock ? »}). Cependant, elle échoue souvent sur des requêtes précises (\textit{« Où est utilisée la variable 'MySpecificVar' ? »}). Deux noms de variables différents peuvent avoir des vecteurs très proches s'ils apparaissent dans des contextes similaires, menant à des hallucinations.

% \subsubsection{Implémentation du Routeur Intelligent}
% Nous avons introduit un composant \textbf{Router} en amont de la récupération. Il s'agit d'un appel léger au LLM qui classifie l'intention de l'utilisateur :
% \begin{enumerate}
%     \item \textbf{Intent "Conceptuel" :} Activation du \textbf{Retriever FAISS} (Recherche sémantique).
%     \item \textbf{Intent "Exact" :} Activation du \textbf{Retriever Grep}. Ce module exécute des expressions régulières (Regex) directement sur le code source pour une précision de 100\%.
% \end{enumerate}
% Cette hybridité permet au système de bénéficier du "meilleur des deux mondes" : la compréhension du sens et la précision chirurgicale de la recherche textuelle.

\subsection{Amélioration progressive de l'architecture grâce à LangGraph}

Plutôt qu'une pipeline linéaire (Question $\rightarrow$ Recherche $\rightarrow$ Réponse), nous avons implémenté une architecture cyclique basée sur des graphes d'états (\texttt{StateGraph}), utilisant la librairie \textbf{LangGraph} \cite{LangGraph}.

\subsubsection{Évolution de la pipeline}
La pipeline prévisionnelle que nous avons établie dans la proposition détaillé était la suivante : 

\begin{figure}[H]
\begin{center}
\begin{tikzpicture}[
    node distance=1cm,
    auto,
    block/.style={
        rectangle,
        draw=cyan!60,
        fill=cyan!20,
        text width=7em,
        text centered,
        rounded corners,
        minimum height=3em
    },
    line/.style={
        draw,
        -latex,
        thick,
        color=black
    }
]

    \node [block] (challenges) {Challenges Q/A pairs};
    \node [block, right=of challenges] (questions) {Questions};
    \node [block, above=of challenges] (parser) {Parser};
    \node [block, above=of parser] (data) {Data Base};
    \node [block, right=of parser] (rag) {Retrieval Augmented Generation (RAG)};
    \node [block, right=of questions] (prompt) {Engineered Prompt};
    \node [block, below=of questions] (ref) {Reference Answers (for supervision)};
    \node [block, below right=of ref, xshift=0.3cm] (grader) {Answer Grader};
    \node [block, below=of grader] (benchmark) {Benchmark};
    \node [block, right=of benchmark] (answer) {Final Answer};
    \node [block, right=of prompt] (llm) {Main LLM};
    \node [block, below=of llm] (logic) {Logic Checker};

    \path [line] (challenges) -- (questions);
    \path [line] (challenges) |- (ref);
    \path [line] (ref) |- (grader);
    \path [line] (logic) |- (grader);
    \draw [-latex] (logic) -- node[pos=0.25, right] {else} (answer);
    \draw [-latex] (logic) -| node[pos=0.3, below] {if error detected} (prompt);
    \path [line] (grader) -- (benchmark);
    \path [line] (questions) -- (prompt);
    \path [line] (questions) -- (rag);
    \path [line] (prompt) -- (llm);
    \path [line] (llm) -- (logic);
    \path [line] (data) -- (parser);
    \path [line] (parser) -- (rag);
    \path [line] (rag) -| (prompt);
    \path [line] (questions) -- (prompt);

\end{tikzpicture}
\end{center}
\caption{Pipeline prévisionnelle issue de la proposition détaillée}
\end{figure}

Dans un premier temps, nous l’avons implémentée sans coder de boucle rétroactive mais nous nous sommes rendu compte que le RAG était peu adapté à certains types de questions comme les questions de comptage. Nous avons donc remplacé le bloc RAG entre Questions et Engineered Prompt par la sous-pipeline suivante :


\begin{figure}[H]
\begin{center}
\begin{tikzpicture}[
    node distance=1cm,
    auto,
    block/.style={
        rectangle,
        draw=cyan!60,
        fill=cyan!20,
        text width=7em,
        text centered,
        rounded corners,
        minimum height=3em
    },
    line/.style={
        draw,
        -latex,
        thick,
        color=black
    }
]

    \node [block] (question) {Question};
    \node [block, right=of question] (router) {Router};
    \node [block, above right=of router] (grep) {Grep Search};
    \node [block, below right=of router] (rag) {Vector Search (RAG Fusion)};
    \node [block, above right=of rag] (docs) {List of Documents};

    \path [line] (question) -- (router);
    \path [line] (router) -- (grep);
    \path [line] (router) -- (rag);
    \path [line] (grep) -- (docs);
    \path [line] (rag) -- (docs);

\end{tikzpicture}
\end{center}
\caption{Workflow avec routeur qui remplaçait le bloc RAG dans la pipeline prévisionnelle}
\end{figure}

Cette nouvelle architecture était plus performante mais manquait d’une boucle rétroactive pour les questions complexes. Or plusieurs études récentes démontrent l'intérêt de passer d'un pipeline RAG statique à un système d'agents autonomes permettant plusieurs itérations de réflexion \cite{Singh2025} \cite{Liang2025}. Nous avons donc repensé la pipeline prévisionnelle pour que cette boucle soit vraiment le cœur du modèle et avons abouti à ce workflow agentique : 

\begin{figure}[H]
\begin{center}
\begin{tikzpicture}[
    node distance=1cm,
    auto,
    block/.style={
        rectangle,
        draw=cyan!60,
        fill=cyan!20,
        text width=7em,
        text centered,
        rounded corners,
        minimum height=3em
    },
    line/.style={
        draw,
        -latex,
        thick,
        color=black
    }
]

    \node [block] (challenges) {Challenges Q/A pairs};
    \node [block, right=of challenges] (questions) {Questions};
    \node [block, right=of questions] (planner) {Planner LLM};
    \node [block, above=of planner] (tool) {Tool (RAG, Grep, ...)};
    \node [block, right=of tool] (distillation) {Distillation LLM};
    \node [block, below=of challenges] (ref) {Reference Answers (for supervision)};
    \node [block, right=of ref] (grader) {Answer Grader};
    \node [block, below=of grader] (benchmark) {Benchmark};
    \node [block, right=of prompt] (llm) {Main LLM};
    \node [block, below=of llm, yshift=-0.185cm] (answer) {Final Answer};

    \path [line] (challenges) -- (questions);
    \path [line] (challenges) -- (ref);
    \path [line] (ref) -- (grader);
    \draw [-latex] (planner) |- node[pos=0.24, right] {else} (answer);
    \draw [-latex] (planner) -- node[pos=0.5, left, align=right] {first pass or \\ unsatisfactory answer} (tool);
    \path [line] (planner) |- (grader);
    \path [line] (grader) -- (benchmark);
    \path [line] (questions) -- (planner);
    \path [line] (tool) -- (distillation);
    \path [line] (distillation) -- (llm);
    \path [line] (llm) -- (planner);

\end{tikzpicture}
\end{center}
\caption{Pipeline Agentique qui correspond à l'état actuel de notre architecture}
\end{figure}

À ce jour, cette architecture reste au cœur de notre système, à une optimisation près. Afin d'éviter de solliciter inutilement le \texttt{Main LLM} lorsqu'un outil ne renvoie aucun résultat (en raison d'arguments invalides ou d'une recherche infructueuse), nous redirigeons immédiatement le retour d'erreur vers le \texttt{Planner LLM}. Cette boucle de correction rapide, appelée \textit{follow-up call} dans notre base de code, oblige le planificateur à ajuster ses paramètres sans répéter son erreur initiale. L'avantage technique de cette approche réside dans l'exploitation du \textit{KV-caching} : le modèle conservant en mémoire les matrices d'attention déjà calculées pour une séquence initiale de tokens identique, cette itération s'appuie directement sur le contexte de la génération précédente. Ainsi, le coût de calcul additionnel est minime, puisque l'attention n'est recalculée que pour le nouveau segment de prompt ajouté à la fin de l'historique.

\subsubsection{Fonctionnement de la boucle}
Le cycle se déroule comme suit :
\begin{enumerate}
    \item \textbf{Planification :} L'agent analyse la demande. S'il lui manque des informations, il décide d'appeler un outil (\texttt{Tool Call}).
    \item \textbf{Exécution :} L'outil (RAG ou Grep) s'exécute et retourne des données brutes.
    \item \textbf{Observation et Réflexion :} L'agent reçoit ces données. Il peut décider qu'elles sont suffisantes et générer la réponse, ou qu'elles sont incomplètes et formuler une nouvelle requête (ex: "J'ai trouvé la table, maintenant je cherche la colonne").
\end{enumerate}

Cette capacité de rétroaction et d'itération est ce qui distingue fondamentalement notre prototype d'un simple moteur de recherche : il mime le comportement d'un développeur humain qui navigue dans le code.

\subsection{Outils}

Pour permettre au modèle d'appeler les outils que l'on a définis, nous utilisons la syntaxe dite "OpenAI-compatible", introduite par OpenAI et prise en charge par les APIs de la plupart des grands modèles. Lorsque l'on veut que le modèle appelle un tool, nous lui fournissons la liste des tools au format suivant :

\begin{verbatim}
{
    "type": "function",
    "function": {
        "name": "tool_name",
        "description": "Ce que fait le tool",
        "parameters": {
            "type": "object",
            "properties": { /* arguments */ },
            "required": [ /* arguments obligatoires */ ]
        }
    }
}
\end{verbatim}

Cette syntaxe force (ou, pour certains modèles, encourage fortement mais sans garantie) le choix de l'un des outils listés. Le modèle répond dans un format standardisé indiquant notamment l'outil et les arguments choisis.

\subsubsection{Script Finder Tool et Gestion Centralisée des Chemins}

Afin de permettre à l'agent de consulter l'intégralité d'un code source lorsqu'une vue globale est nécessaire, un premier outil que nous avons implémenté est le \texttt{script\_finder\_tool}. Ce développement a également été l'occasion de concevoir une gestion centralisée de l'arborescence des fichiers utilisée par tous les outils. 

En pratique, les scripts extraits des comptes clients se présentent souvent sous la forme d'un \textit{dump} de fichiers plats nommés par des identifiants bruts (par exemple \texttt{67984.nvn}). Pour que l'agent puisse raisonner sur une structure de projet cohérente, notre système reconstruit une arborescence virtuelle complète à l'aide d'un simple fichier de correspondance (\texttt{mapping.txt}). Cette architecture rend notre solution extrêmement modulaire : pour analyser un nouveau compte client, il suffit de fournir le dossier contenant les scripts en vrac et son fichier de \textit{mapping} associé. Le système gère alors automatiquement et de manière transparente la traduction entre les identifiants techniques et les chemins métiers compréhensibles par le LLM.

Le \texttt{script\_finder\_tool} s'appuie sur cette mécanique pour retrouver et lire n'importe quel fichier à partir d'un fragment de son chemin d'origine. Cependant, dans sa version actuelle, cet outil renvoie systématiquement l'intégralité du texte du script demandé. En raison de la saturation rapide de la fenêtre de contexte que cela engendre, nous avons dû brider son utilisation via les directives système (\textit{« Use RARELY and only when necessary due to high token cost »}), forçant ainsi le \texttt{Planner LLM} à privilégier des outils plus ciblés comme le RAG ou le grep.

Une piste d'amélioration solide consisterait à faire évoluer cet outil pour qu'il accepte des arguments de plage de lecture (par exemple \texttt{start\_line} et \texttt{end\_line}). Cela permettrait à l'agent d'inspecter chirurgicalement le code situé juste au-dessus ou en dessous d'un \textit{chunk} renvoyé par le RAG, sans inonder sa mémoire de travail. Si l'implémentation Python de cette fonctionnalité est triviale, nous avons sciemment choisi de ne pas l'intégrer au cours des dernières semaines du PSC. En effet, dans un flux de travail agentique, atteindre un point d'équilibre où le modèle choisit judicieusement entre plusieurs outils concurrents requiert un \textit{prompt engineering} millimétré. Modifier l'utilisation du \texttt{script\_finder\_tool} risquait de perturber celle du RAG et donc de déstabiliser cet équilibre précaire que nous avions atteint. 

Néanmoins, l'intégration de cette pagination ciblée reste une évolution naturelle et nécessaire si l'on souhaite se rapprocher des standards actuels du marché des agents codeurs autonomes (comme GitHub Copilot ou Cursor).

\subsubsection{Tree tool}
Pour fournir l'arborescence de certains fichiers sources au modèle, nous lui donnons un outil lui permettant de recevoir la représentation de l'arborescence au format usuel employé par exemple par l'outil tree de Linux, et donc plus facilement compris par le modèle.

En pratique, cet outil est rarement appelé explicitement par le planificateur car nous donnons une version condensée de l'arborescence de la base de code au début de chaque prompt.

\subsubsection{RAG tool: l'outil de recherche sémantique}

L’implémentation initiale de notre système RAG reposait sur FAISS (Facebook AI Similarity Search), une solution efficace pour le prototypage rapide, mais fondamentalement limitée par son approche unidimensionnelle. Dans cette première version, la recherche s'effectuait exclusivement via des vecteurs denses (\textit{embeddings}) qui capturaient le sens sémantique global des requêtes, mais peinaient face au vocabulaire très spécifique d'une base de code (noms de variables exacts, UUIDs, ou acronymes métiers de Lokad/Envision). Pour pallier cette faiblesse et permettre à l'agent de cibler des fichiers spécifiques ou des mots-clés, nous avions mis en place un système de reclassement manuel codé en dur en Python. Après avoir récupéré les \textit{K} meilleurs résultats de FAISS, le script appliquait des multiplicateurs mathématiques artificiels : le score d'un chunk était augmenté de quelques pourcents s'il contenait un mot-clé demandé, ou pénalisé s'il ne correspondait pas à l'expression régulière (\textit{regex}) du chemin de fichier spécifié. Bien que fonctionnelle sur une petite échelle, cette méthode était algorithmiquement fragile. Elle obligeait le système à manipuler des scores post-recherche, ce qui signifiait que si FAISS ne remontait pas le bon fichier dans son top 30 initial, le système de pénalités Python ne pouvait tout simplement pas corriger le tir, la bonne réponse ayant déjà été éliminée de la sélection.\\

Pour franchir un cap en termes de précision, nous avons migré la base vectorielle vers Qdrant, ce qui nous a permis de débloquer la puissance de la recherche hybride en combinant une base de données "Dense" (comme FAISS) et une base de données "Sparse" qui permet de faire des recherches par fréquence de tokens. Pour cette base de données sparse, nous utilisons BM25 (Best Matching 25) \cite{Robertson2009} qui est une fonction de classement fréquentielle estimant la pertinence d'un document $D$ par rapport à une requête $Q$. Le score de pertinence est calculé selon la formule suivante :

\begin{equation}
\text{Score}(D, Q) = \sum_{q \in Q} \text{IDF}(q) \cdot \frac{f(q, D) \cdot (k_1 + 1)}{f(q, D) + k_1 \cdot (1 - b + b \cdot \frac{|D|}{\text{avgdl}})}
\end{equation}

où $q\in Q$ sont les tokens compris dans $Q$. Pour bien appréhender la mécanique de cette formule, il convient de la décomposer en trois axes fondamentaux qui agissent de concert : \\

\begin{itemize}
    \item \textbf{Le poids intrinsèque du terme (La composante IDF) :} 
    La partie $\text{IDF}(q)$ agit comme un filtre de rareté. Elle est classiquement calculée selon la formule suivante :
    $$\text{IDF}(q) = \ln\left(\frac{N - n(q) + 0.5}{n(q) + 0.5} + 1\right)$$
    où $N$ représente le nombre total de documents dans la collection (la base de code), et $n(q)$ est le nombre de documents contenant le terme $q$. 
    Elle part du principe mathématique qu'un mot très commun dans la base de code (où $n(q)$ s'approche de $N$) verra son score s'effondrer vers zéro, car il apporte peu d'information discriminante. À l'inverse, un terme technique rare (où $n(q)$ est très petit) verra son score multiplié de façon significative.\\
    
    \item \textbf{La saturation de la fréquence (Le paramètre $k_1 \in [1.2, 2.0]$) :} 
    La fraction centrale de l'équation gère la fréquence d'apparition du terme $f(q, D)$, c'est-à-dire le nombre de fois où le mot est trouvé dans le document. Sans le paramètre constant $k_1$, un document répétant 100 fois le même mot écraserait tous les autres résultats. Ce paramètre $k_1$ impose une courbe asymptotique : la première apparition d'un mot donne un score fort, la deuxième un peu moins, jusqu'à plafonner mathématiquement vers $(k_1 + 1)$. Cela traduit l'idée sémantique que "mentionner un sujet 5 fois prouve qu'on en parle, le mentionner 50 fois n'apporte plus de preuve supplémentaire".\\

    \item \textbf{La normalisation par la longueur (Le paramètre $b \approx 0.75$) :} 
    Le dénominateur intègre le ratio $\frac{|D|}{\text{avgdl}}$, qui compare la longueur du document évalué $|D|$ à la longueur moyenne des documents de la base \texttt{avgdl}. Par pure probabilité, un fichier de 3000 lignes a beaucoup plus de chances de contenir accidentellement les mots de la requête qu'un petit fichier de 50 lignes. Le paramètre constant $b$ (généralement fixé à $0.75$, mais obligatoirement compris entre $0$ et $1$) vient pondérer cette injustice statistique en pénalisant les documents excessivement longs, assurant ainsi que la concision et la densité de l'information sont récompensées.\\
\end{itemize}

Avec Qdrant, nous avons mis en place une architecture sophistiquée à 4 flux de recherche simultanés (2 flux \textit{denses} et 2 flux \textit{sparses}) qui se balancent mutuellement grâce au Reciprocal Rank Fusion (RRF) \cite{Cormack2009}, qui favorise les documents récupérés dans plusieurs flux différents :
$$RRF(d)=\sum_{r\in R}\frac{1}{k+r(d)}$$
où chaque élément $r\in R$ correspond à un flux et associe à chaque document $d$ son classement pour celui-ci, et $k$ est une constante qui est usuellement choisie égale à 60.

Les deux premiers flux opèrent au niveau global sur toute la base de code : le flux \textit{Dense} cherche les concepts, tandis que le flux \textit{Sparse} (BM25) vérifie que les mots demandés sont présents. Les deux flux suivants sont conditionnels et dédiés au "Source Boosting" : si l'agent spécifie un chemin de fichier, Qdrant relance une recherche Dense et une recherche Sparse uniquement à l'intérieur de ces fichiers cibles. Cette architecture matricielle (cf. Table \ref{tab:matrice_flux}) garantit que peu importe si l'utilisateur pose une question large et conceptuelle ou s'il cherche l'usage exact d'une fonction spécifique dans un dossier précis, le système a mathématiquement la capacité de trouver le bon chunk sans nécessiter de calcul post-recherche.\\

\begin{table}[htbp]
\centering
\renewcommand{\arraystretch}{1.5}
\begin{tabularx}{\textwidth}{l >{\raggedright\arraybackslash}X >{\raggedright\arraybackslash}X}
\toprule
\textbf{Portée / Méthode} & \textbf{Dense (Sémantique)} & \textbf{Sparse (Lexical / BM25)} \\ 
\midrule
\textbf{Global} & 
\textbf{Flux 1 :} Capture les concepts, la logique métier et les intentions globales. & 
\textbf{Flux 2 :} Capture la syntaxe exacte, les noms de variables et les identifiants techniques. \\ 
\addlinespace
\textbf{Filtré (Sources)} & 
\textbf{Flux 3 (Boosté) :} Recherche conceptuelle limitée aux fichiers cibles. \newline \textit{Seuil : 0.3} & 
\textbf{Flux 4 (Boosté) :} Recherche de mots-clés exacte limitée aux fichiers cibles. \newline \textit{Seuil : 2.0} \\ 
\bottomrule
\end{tabularx}
\caption{Matrice de recherche hybride fusionnée via Reciprocal Rank Fusion (RRF).}
\label{tab:matrice_flux}
\end{table}

Le principal défaut de cette architecture RAG est que le RRF est aveugle aux scores réels et ne fusionne que les classements (rangs). Si l'agent LLM hallucine et demande de chercher le mot-clé \texttt{forecasting} dans un fichier source qui n'a rien à voir (ex: \texttt{/1. utilities/Modules/Colors}), les flux 3 et 4 trouveraient quand même des résultats dans ce fichier, et le rang \#1 de ces flux obtiendrait un boost massif, enterrant la vraie réponse. Pour contrer partiellement cela, nous avons équipé Qdrant de seuils de score pour les flux de source boosting et nous les avons ajustés expérimentalement. Ainsi, si la meilleure correspondance dans le fichier cible n'atteint pas un seuil minimal de pertinence (cosine-similarity ou BM25), la recherche ciblée est purement et simplement avortée. Le système RRF retombe alors automatiquement sur les flux globaux (1 et 2), sauvant la recherche et garantissant que le boost de source n'est appliqué que lorsqu'il est factuellement mérité, offrant une meilleure résilience face aux hallucinations de l'agent. \\

Enfin, nous avons ajouté une ultime couche de raffinement contextuel : le Reranking avec un Cross-Encoder (Jina AI). Qdrant agit comme un filet large et intelligent, remontant les 15 ou 30 meilleurs candidats, mais ces modèles utilisent la similarité cosinus, qui reste une comparaison mathématique relativement simple. Le modèle \texttt{jinaai/jina-reranker-v2-base-multilingual} prend le relais en utilisant un \textit{transformer} pour lire simultanément la question de l'utilisateur et le code récupéré. Nous avons choisi Jina spécifiquement pour son excellence dans un environnement bilingue, lui permettant de comprendre parfaitement le lien sémantique entre une question posée en français, des chemins de fichiers français (\texttt{/4. Optimization workflow/French/...}), et de la logique métier codée en anglais. Ce modèle ne fait pas que vérifier la présence des mots ; il comprend profondément le contexte technique et évalue si le chunk répond réellement à la question.\\

Dans le but de maximiser l'efficacité du reranking, nous ne passons pas simplement le code brut au Cross-Encoder, nous concaténons le chemin du fichier au sommet du document (\texttt{Source: ...}), et nous structurons la requête en y intégrant explicitement les directives de l'agent (\texttt{Required keywords: ... Target source: ...}). Le transformeur Jina peut ainsi lier son attention directement entre le chemin demandé et la métadonnée du chunk. Les scores bruts (logits) générés par Jina sont ensuite normalisés via une fonction Sigmoïde pour obtenir un score normalisé entre 0.0 et 1.0. Cette architecture finale combine la vitesse et la couverture exhaustive de la base Qdrant avec l'analyse chirurgicale et sémantique d'un modèle neuronal lourd, garantissant que seuls les extraits de code les plus pertinents et exacts atterrissent dans la fenêtre de contexte de l'agent final.
 
\subsubsection{Grep tool : la recherche lexicale exacte guidée par la structure du code}

Le \texttt{grep\_tool} répond à un besoin complémentaire du \textit{RAG}. Lorsqu'une question porte sur un identifiant précis, un chemin exact, un littéral ou une construction syntaxique bien déterminée, une recherche sémantique pure devient trop floue. Dans ce cas, nous préférons un outil de recherche textuelle exacte, piloté par expressions régulières.

Cependant, notre \texttt{grep\_tool} n'est pas un simple appel à \texttt{rg} ou à \texttt{grep} sur des fichiers bruts. Il opère sur notre représentation déjà parsée du code Envision : les scripts sont d'abord découpés en blocs syntaxiques cohérents (\texttt{IMPORT}, \texttt{READ}, \texttt{WRITE}, \texttt{SHOW}, \texttt{EXPORT}, etc.), puis la recherche s'applique sur ces blocs. Cela permet de restreindre la recherche à certaines formes de code, ce qui est particulièrement utile pour répondre à des questions comme \textit{« où lit-on \code{/Clean/Items.ion} ? »} ou \textit{« dans quels inspecteurs apparaît un \texttt{show table} ? »}.

\begin{figure}[H]
\begin{center}
\resizebox{0.97\textwidth}{!}{
\begin{tikzpicture}[
    >=Latex,
    box/.style={
        rectangle,
        rounded corners,
        draw=cyan!60!black,
        fill=cyan!12,
        align=center,
        minimum height=1.35cm,
        text width=3.8cm,
        inner sep=6pt
    },
    box2/.style={
        rectangle,
        rounded corners,
        draw=orange!70!black,
        fill=orange!15,
        align=center,
        minimum height=1.35cm,
        text width=4.2cm,
        inner sep=6pt
    },
    box3/.style={
        rectangle,
        rounded corners,
        draw=green!50!black,
        fill=green!12,
        align=center,
        minimum height=1.35cm,
        text width=4.0cm,
        inner sep=6pt
    },
    line/.style={draw, thick, -Latex}
]
    \node[box] (scripts) at (0,0) {Scripts miroirs\\\code{env_scripts/*.nvn}};
    \node[box2] (parser) at (4.9,0) {\texttt{EnvisionParser}\\blocs \texttt{IMPORT}, \texttt{READ}, \texttt{SHOW}, \dots};
    \node[box2] (query) at (4.9,-2.5) {\texttt{grep\_tool}\\\texttt{pattern} + \texttt{sources} + \texttt{block\_type}};
    \node[box3] (matches) at (9.8,-2.5) {Blocs correspondants\\et scripts sources};
    \node[box3] (compact) at (14.7,-2.5) {Compaction contextuelle\\fenêtre de lignes utiles};

    \draw[line] (scripts) -- (parser);
    \draw[line] (parser) -- (query);
    \draw[line] (query) -- (matches);
    \draw[line] (matches) -- (compact);
\end{tikzpicture}}
\end{center}
\caption{Chaîne de traitement du \texttt{grep\_tool} : la recherche ne s'effectue pas sur du texte brut, mais sur des blocs Envision parsés puis compactés avant renvoi au modèle.}
\end{figure}

\paragraph{Trois leviers de contrôle}
L'outil est construit autour de trois paramètres simples mais très efficaces :
\begin{itemize}
    \item \texttt{pattern} : le motif recherché, généralement un identifiant, une chaîne ou une petite expression régulière.
    \item \texttt{sources} : un filtre optionnel sur les chemins de scripts, pour restreindre la recherche à une zone du dépôt.
    \item \texttt{block\_type} : une liste optionnelle de types de blocs Envision, permettant par exemple de ne chercher que dans les \texttt{READ} ou les \texttt{SHOW}.
\end{itemize}

Cette combinaison donne au planificateur un contrôle beaucoup plus fin qu'une recherche globale. Il peut d'abord cibler un dossier logique, puis resserrer encore la recherche au niveau d'une structure syntaxique particulière.

\begin{figure}[H]
\begin{center}
\resizebox{0.96\textwidth}{!}{
\begin{tikzpicture}[
    >=Latex,
    knob/.style={
        rectangle,
        rounded corners,
        draw=blue!60!black,
        fill=blue!10,
        align=center,
        minimum height=1.2cm,
        text width=4cm,
        inner sep=6pt
    },
    result/.style={
        rectangle,
        rounded corners,
        draw=green!50!black,
        fill=green!12,
        align=center,
        minimum height=1.25cm,
        text width=5.2cm,
        inner sep=6pt
    },
    line/.style={draw, thick, -Latex}
]
    \node[knob] (pattern) at (0,0) {\texttt{pattern}\\\code{StockEvol}\\ou \code{/Clean/Items.ion}};
    \node[knob] (sources) at (5.2,0) {\texttt{sources}\\\code{/Modules/Functions}\\ou \code{/3. Inspectors/}};
    \node[knob] (blocks) at (10.4,0) {\texttt{block\_type}\\\texttt{[READ, FORM\_READ]}\\ou \texttt{[SHOW]}};
    \node[result] (target) at (5.2,-2.6) {Recherche exacte mais ciblée\\sur une zone et une structure précise du code};

    \draw[line] (pattern) -- (target);
    \draw[line] (sources) -- (target);
    \draw[line] (blocks) -- (target);
\end{tikzpicture}}
\end{center}
\caption{Les trois degrés de liberté du \texttt{grep\_tool}. Le planificateur peut jouer sur le motif, la zone du dépôt et le type syntaxique recherché.}
\end{figure}

Une difficulté majeure que nous avons dû surmonter concerne la définition itérative des noms de fichiers dans Envision. Le code utilise fréquemment des interpolations de chaînes de caractères pour définir les chemins d'importation, par exemple :
\begin{verbatim}
const testStorage = "/Data/"
const inputPath = "\{testStorage}Clean/"
read "\{inputPath}Items.ion" as Items
\end{verbatim}
Une recherche textuelle simple (Grep) sur la chaîne \texttt{"/Clean/Items.ion"} échoue ici, car le chemin complet n'est jamais écrit explicitement en une seule ligne.
\\
Nous avons donc dû ajuster l'outil Grep pour gérer cette complexité. L'outil effectue désormais une résolution récursive à la volée : il isole la dernière partie du chemin, vérifie si le préfixe correspond à une variable définie précédemment (comme \texttt{\{inputPath\}}), et remonte ainsi la chaîne de dépendances jusqu'à la racine. Bien que fonctionnelle, cette méthode est coûteuse en temps de calcul car elle s'exécute au moment de la requête (\textit{runtime}).

\paragraph{Exemples réels sur le dépôt courant}
Les requêtes ci-dessous ont été exécutées sur l'état actuel du dépôt.
\begin{itemize}
    \item La requête \code{pattern="StockEvol"} avec \code{sources="/Modules/Functions"} retourne un unique résultat, le module \code{/1. utilities/Modules/Functions.nvn}. Ce cas est typique d'une recherche d'identifiant précis.
    \item La requête \code{pattern="show table"}, \code{sources="/3. Inspectors/"} et \code{block_type=[SHOW]} retourne 16 résultats, tous situés dans les inspecteurs. Ici, le filtre syntaxique évite de rechercher inutilement dans les blocs de calcul.
    \item La requête \code{pattern="/Clean/Items.ion"} avec \code{block_type=[READ, FORM_READ]} retourne 31 résultats. Le \texttt{grep\_tool} identifie alors précisément les scripts qui lisent ce fichier, avant de compacter les extraits pertinents.
\end{itemize}

\begin{verbatim}
read "/Clean/Items.ion" as Items[Id unsafe] with
  "Sku" as Id : text
  Name : text
  Category : text
  SellPrice : number
  BuyPrice : number
  StockOnHand : number
  ...
\end{verbatim}

L'extrait ci-dessus illustre le format final transmis au modèle : le bloc est raccourci autour de la zone utile, avec une fenêtre de contexte ajustée automatiquement pour respecter un budget total de lignes.

\paragraph{Apport pour l'agent}
Le \texttt{grep\_tool} joue donc le rôle d'un instrument de \textbf{précision lexicale}. Là où le \textit{RAG} répond mieux aux formulations conceptuelles et où le \texttt{graph\_tool} fournit des relations structurelles, le \texttt{grep\_tool} excelle pour :
\begin{itemize}
    \item retrouver un identifiant exact ;
    \item confirmer qu'un chemin ou un littéral apparaît réellement dans le code ;
    \item restreindre une recherche à un dossier particulier ;
    \item cibler une forme syntaxique précise, comme les \texttt{READ} ou les \texttt{SHOW}. \\
\end{itemize}

Cette complémentarité est importante dans notre architecture : elle permet au planificateur de choisir entre une recherche \textbf{sémantique}, \textbf{lexicale} ou \textbf{structurelle} selon la nature effective de la question.


\subsubsection{Graph tool : un graphe de dépendances structurelles}

Le \texttt{graph\_tool} complète le \textit{RAG} sémantique et la recherche lexicale par une troisième vue du code : une vue \textbf{structurelle}. L'idée n'est plus seulement de retrouver des morceaux de texte pertinents, mais de reconstruire explicitement les dépendances du projet sous la forme d'un graphe navigable.

À la différence du graphe d'orchestration \textit{LangGraph}, qui pilote la boucle de raisonnement de l'agent, le \texttt{graph\_tool} représente le contenu même du dépôt Envision. Chaque script, chaque dossier logique, chaque fichier de données, ainsi que certaines entités internes (tables, fonctions) deviennent des nœuds ; les relations \texttt{imports}, \texttt{reads}, \texttt{writes}, \texttt{defines}, \texttt{contains} et \texttt{sibling} deviennent des arêtes. Ce graphe permet donc de répondre à des questions structurelles du type : \textit{« quels scripts lisent \code{/Clean/Items.ion} ? »}, \textit{« quel module est importé par tel inspecteur ? »}, ou encore \textit{« dans quel script une table est-elle définie ? »}.

\paragraph{Deux arbres fondamentaux reliés}
Le graphe n'est pas un arbre unique. Il juxtapose deux hiérarchies distinctes :
\begin{itemize}
    \item le domaine \textbf{scripts}, qui organise les fichiers \texttt{.nvn} selon leur arborescence logique ;
    \item le domaine \textbf{data}, qui organise les fichiers \texttt{.ion}, \texttt{.csv} ou \texttt{.xlsx} manipulés par les scripts.
\end{itemize}
Ces deux arbres sont ensuite reliés par des arêtes transverses de lecture et d'écriture. Le graphe encode ainsi en une seule structure à la fois la hiérarchie du code et le flot des données.

\begin{figure}[H]
\begin{center}
\resizebox{0.96\textwidth}{!}{
\begin{tikzpicture}[ 
    >=Latex, 
    script/.style={ 
        rectangle, 
        rounded corners, 
        draw=blue!60!black, 
        fill=blue!10, 
        minimum height=0.95cm, 
        text width=4.4cm, 
        align=center, 
        inner sep=5pt 
    }, 
    data/.style={ 
        rectangle, 
        rounded corners, 
        draw=green!50!black, 
        fill=green!10, 
        minimum height=0.95cm, 
        text width=4.2cm, 
        align=center, 
        inner sep=5pt 
    }, 
    line/.style={draw, thick, -Latex}, 
    % Styles mis à jour : "text=" assure que le texte prend la même couleur
    relImport/.style={draw=orange!85!black, text=orange!85!black, thick, -Latex},
    relRead/.style={draw=red!70!black, text=red!70!black, dashed, thick, -Latex}
] 
    % --- NOEUDS SCRIPTS ---
    \node[script] (sroot) at (0,4.7) {\textbf{scripts : /}}; 
    
    \node[script] (s1) at (-5.2,3.0) {\code{/1. utilities}}; 
    \node[script] (s3) at (0,3.0) {\code{/3. Inspectors}}; 
    \node[script] (s4) at (5.2,3.0) {\shortstack{\code{/4. Optimization}\\\code{workflow}}}; 
    
    \node[script] (s1_mod) at (-5.2,1.3) {\shortstack{\code{/1. utilities/Modules/}\\\code{Global Parameters}}}; 
    \node[script] (s3_sales) at (0,1.3) {\shortstack{\code{/3. Inspectors/}\\\code{2 - Sales Analysis.nvn}}}; 

    % --- NOEUDS DATA ---
    \node[data] (droot) at (0,-1.0) {\textbf{data : /}}; 
    \node[data] (din) at (-5.2,-2.7) {\code{/Input}}; 
    \node[data] (dclean) at (0,-2.7) {\code{/Clean}}; 
    \node[data] (dman) at (5.2,-2.7) {\code{/Manual}}; 
    
    \node[data] (items) at (-5.2,-4.4) {\code{/Clean/Items.ion}}; 
    \node[data] (orders) at (0,-4.4) {\code{/Clean/Orders.ion}}; 
    \node[data] (stock) at (5.2,-4.4) {\code{/Clean/Stock.ion}}; 

    % --- LIENS STANDARDS (Noirs) ---
    \draw[line] (sroot) -- (s1); 
    \draw[line] (sroot) -- (s3); 
    \draw[line] (sroot) -- (s4); 
    
    \draw[line] (s1) -- (s1_mod);
    \draw[line] (s3) -- (s3_sales);

    \draw[line] (droot) -- (din); 
    \draw[line] (droot) -- (dclean); 
    \draw[line] (droot) -- (dman); 
    
    \draw[line] (dclean) -- (items); 
    \draw[line] (dclean) -- (orders); 
    \draw[line] (dclean) -- (stock); 

    % --- LIENS DE RELATIONS (Couleurs) ---
    % Imports (Orange)
    \draw[relImport] (s3_sales.west) -- node[above] {imports} (s1_mod.east); 
    
    % Reads (Rouges pointillés)
    \draw[relRead] (s3_sales.south west) -- node[left] {reads} (items.north); 
    \draw[relRead] (s3_sales.south east) -- node[right] {reads} (stock.north); 
    
    % Contournement fluide par la droite pour la flèche du milieu
    \draw[relRead] (s3_sales.south) .. controls +(2.6, -1.5) and +(2.6, 1.5) .. node[right] {reads} (orders.north); 
\end{tikzpicture}}
\end{center}
\caption{Visualisation conceptuelle du graphe interne : deux arbres fondamentaux (\textit{scripts} et \textit{data}) reliés par des arêtes de dépendance.}
\end{figure}

\paragraph{Exemple local de sous-graphe}
Une fois le graphe construit, l'agent peut raisonner localement autour d'un nœud particulier. La figure suivante montre le sous-graphe centré sur \code{/3. Inspectors/2 - Sales Analysis.nvn}, qui importe un module, lit plusieurs fichiers de données et définit une table.

\begin{figure}[H]
\begin{center}
\resizebox{0.92\textwidth}{!}{
\begin{tikzpicture}[
    >=Latex,
    node distance=1.6cm and 2.1cm,
    center/.style={
        shape=ellipse,
        draw=cyan!70!black,
        fill=cyan!18,
        minimum width=5.6cm,
        minimum height=1.55cm,
        text width=5.2cm,
        align=center,
        inner sep=6pt
    },
    script/.style={
        rectangle,
        rounded corners,
        draw=blue!60!black,
        fill=blue!10,
        minimum height=1cm,
        text width=4cm,
        align=center,
        inner sep=5pt
    },
    data/.style={
        rectangle,
        rounded corners,
        draw=green!50!black,
        fill=green!10,
        minimum height=1cm,
        text width=3.2cm,
        align=center,
        inner sep=5pt
    },
    internal/.style={
        rectangle,
        rounded corners,
        draw=orange!70!black,
        fill=orange!12,
        minimum height=1cm,
        text width=4cm,
        align=center,
        inner sep=5pt
    },
    line/.style={draw, thick, -Latex}
]
    \node[center] (sales) {\code{/3. Inspectors/2} \\ \code{- Sales Analysis.nvn}};
    \node[script, above left=1.8cm and 2.5cm of sales] (module) {\code{/1. utilities/Modules/} \\ {Global Parameters}};
    \node[data, right=3.1cm of sales, yshift=2.0cm] (items) {\code{/Clean/Items.ion}};
    \node[data, right=3.1cm of sales, yshift=0.7cm] (orders) {\code{/Clean/Orders.ion}};
    \node[data, right=3.1cm of sales, yshift=-0.7cm] (stock) {\code{/Clean/Stock.ion}};
    \node[data, right=3.1cm of sales, yshift=-2.0cm] (po) {\code{/Clean/PO.ion}};
    \node[internal, below left=1.8cm and 2.5cm of sales] (table) {\code{...::table::ItemsWeek}};

    \draw[line] (sales) -- node[pos=0.52, above, xshift=0.3cm]{\texttt{imports}} (module);
    \draw[line] (sales.east) -- node[pos=0.58, above, yshift=0.2cm]{\texttt{reads}} (items.west);
    \draw[line] (sales.east) -- node[pos=0.58, above]{\texttt{reads}} (orders.west);
    \draw[line] (sales.east) -- node[pos=0.58, below]{\texttt{reads}} (stock.west);
    \draw[line] (sales.east) -- node[pos=0.58, below, yshift=-0.2cm]{\texttt{reads}} (po.west);
    \draw[line] (sales) -- node[pos=0.52, below, xshift=0.4cm]{\texttt{defines}} (table);
\end{tikzpicture}}
\end{center}
\caption{Sous-graphe typique autour d'un script : le graphe fournit directement les relations structurelles utiles à l'agent.}
\end{figure}

\paragraph{Statistiques réelles du graphe}
Sur l'état courant du dépôt, le graphe généré le 11 avril 2026 contient 300 nœuds et 781 arêtes. La répartition est la suivante : 64 scripts, 73 fichiers de données, 96 tables, 17 fonctions et 50 dossiers. Côté relations, on compte notamment 55 arêtes \texttt{imports}, 270 arêtes \texttt{reads}, 63 arêtes \texttt{writes}, 113 arêtes \texttt{defines} et 185 arêtes \texttt{contains}. Cette densité montre bien que le \textit{graph tool} n'est pas un simple affichage d'arborescence, mais une représentation réellement relationnelle du dépôt.

\paragraph{Usage pratique}
Le graphe peut être interrogé indépendamment de l'agent via le CLI \code{python -m env_graph.network}. En pratique, ce CLI expose les mêmes familles d'opérations que l'outil agentique : lecture des statistiques globales, exploration d'un arbre \texttt{scripts} ou \texttt{data}, recherche d'un nœud par son nom, inspection de ses métadonnées, lecture locale de son contenu et navigation dans ses voisins entrants ou sortants. Il permet donc aussi bien une exploration manuelle du dépôt qu'une validation rapide de ce que l'agent utilise ensuite dans la boucle de raisonnement.

\paragraph{Apport pour l'agent}
Le \texttt{graph\_tool} apporte donc une forme de \textbf{mémoire structurelle précompilée}. Là où le \textit{RAG} sémantique répond bien aux questions de logique métier et où le \texttt{grep} répond bien aux recherches exactes, le graphe fournit une vérité relationnelle sur le dépôt :
\begin{itemize}
    \item qui lit quoi ;
    \item qui écrit quoi ;
    \item qui importe quoi ;
    \item dans quel dossier logique se trouve chaque élément ;
    \item quelles tables ou fonctions sont définies par quel script.
\end{itemize}

Cette couche est particulièrement précieuse dans Envision car de nombreux chemins sont construits dynamiquement par interpolation de chaînes. Le graphe capture et normalise une partie de ces dépendances en amont, ce qui réduit la charge d'inférence précédemment évoquée avec le \texttt{grep\_tool}.

\subsection{Query Transformers}

Le problème que l'on peut rencontrer si l'on procède à un simple embedding de la question de l'utilisateur avant \textit{retrieval} dans l'index est l'écart sémantique entre question et documents indexés. Les query transformers cherchent à combler cet écart en retraitant la question par reformulation ou enrichissement.

\subsubsection{RAG Fusion}
Le RAG fusion \cite{RAGFusion} consiste à d’abord demander à un LLM de reformuler la question en plusieurs sous-questions, puis utiliser le RAG pour obtenir les documents les plus adaptés à chacune de ces questions. On classe ensuite les résultats à l’aide de la formule de RRF détaillée plus haut.

En pratique, nous avons constaté que séparer la question en sous-questions fonctionnait moins bien que de proposer plusieurs formulations de la question ayant toutes le même sens. En effet, les réponses aux sous-questions ne peuvent souvent pas être combinées pour répondre à la question initiale alors que les différentes formulations permettent d'élargir le spectre lexical pour la recherche dans les documents. Il est important d'inclure la question d'origine dans la recherche, sans quoi un aspect de la question risque de ne pas se retrouver dans les autres formulations et donc d'être ignoré.

\subsubsection{HyDE}
 La démarche avec HyDE \cite{Jostmann2024} (Hypothetical Document Embedding) est différente. Dans les approches prédécentes, nous cherchions les documents contenant les réponses via l'embedding des questions posées. Le problème avec ce type d'approche était l'écart sémantique qui pouvait exister entre une question et sa réponse. Pour pallier cela, nous générons des documents hypothétiques censés répondre à notre question. Nous utilisons ensuite l'embedding de ces documents idéaux pour le retrieval dans l'index au lieu de l'embedding d'une ou plusieurs questions.
 Un exemple concret : pour la question "quels sont les symptômes de la maladie X ?", il parait intéressant de se concentrer sur les documents qui font intervenir les embeddings de mots comme "fièvre", "fatigue". Sans HyDE, nous ne ciblons pas directement ces mots car nous utilisons l'embedding de la question qui n'inclut que la notion vague de "symptôme".